\section{Control Plane Design}
\label{sec:control-plane}

The control plane is a Python Lambda behind API Gateway. Its job is deliberately narrow: authenticate, derive context, gate modules, minimize surface, and log enough to make incidents explainable. The implementation is small enough to review, but important enough to treat like production security infrastructure.

\subsection{Middleware spine}

Every handler passes through the same middleware chain. In the deployed system, this is implemented with AWS Lambda Powertools decorators \cite{aws_powertools}. The spine has three layers.

\begin{enumerate}
  \item \textbf{Authentication and context.} The middleware validates Cognito JWTs \cite{aws_cognito_user_pools}, resolves tenant membership from the platform table instead of trusting token claims alone, derives role and \scopeprefix{}, and builds the set of available impersonations.
  \item \textbf{Response hardening.} Standard security headers are set at the boundary: HSTS, Content-Security-Policy, \code{X-Frame-Options}, and \code{X-Content-Type-Options}.
  \item \textbf{Conditional revalidation.} The middleware compares the response fingerprint with \code{If-None-Match} and returns \code{304 Not Modified} when the identity contract has not changed.
\end{enumerate}

This is not glamorous code. That is the point. Shared governance should be reviewable and dull.

\subsection{Scope injection}

Tenant scope is derived from the tenant record and injected into downstream calls. It is never accepted from browser input, query parameters, or ad hoc service headers.

For query and analytics modules, the control plane passes \scopeprefix{} as a required field in the internal request payload. Downstream Lambdas reject missing or empty scope with a bad-request error. They re-validate syntax before using the value. They do not independently decide whether a syntactically valid prefix belongs to the caller; that authority remains in the control plane.

This distinction matters. Downstream checks catch malformed or absent scope. They cannot catch a wrong-but-well-formed scope derived by the control plane. The architecture therefore makes the control plane the load-bearing derivation point and makes the residual risk explicit.

\subsection{Module gating and surface minimization}

Each route carries a module requirement. In the Python services this is expressed as a decorator such as \code{@requires_module("grid")}. Before a handler body executes, middleware compares the required module with the enabled-module list attached to the effective tenant context. Missing module, no handler execution.

The same enabled-module set filters the generated OpenAPI document. The \code{/v1/openapi.json} endpoint removes unlicensed paths and prunes unreferenced schema definitions before serving the spec. That avoids a common platform smell: endpoints that are officially documented, operationally denied, and still somehow treated as customer-facing promises.

Filtering happens after impersonation resolution. A support engineer viewing as \code{acme/member} sees the API surface for \code{acme/member}, not the support engineer's broader internal surface.

\subsection{Auditable \viewas{}}

Support engineers sometimes need to see what a customer sees. Shared admin accounts are the tempting shortcut. They are also audit poison. They erase actor attribution, concentrate risk into a single credential, and make every subsequent log line weaker.

The \viewas{} mechanism keeps actor and effective context separate. The actor field remains the authenticated support engineer. The effective context changes to the selected tenant and role. Every audited action can carry both identities.

The web tier signs the \code{X-View-As} payload with HMAC-SHA256. The signed string includes tenant, role, and a per-activation nonce. The signing key is held server-side. The browser never receives it. The nonce is registered in DynamoDB at activation time, and logout or tenant switching revokes it. If nonce registration fails, impersonation fails closed.

A simplified sketch follows. Internal variable names are intentionally omitted.

\begin{lstlisting}[style=python, caption={Simplified view-as signing flow.}]
def activate_view_as(actor, tenant_id, role):
    allowlist = get_available_impersonations(actor)
    if (tenant_id, role) not in allowlist:
        raise Forbidden()
    nonce = create_nonce()
    payload = f"tenant_id={tenant_id};role={role};nonce={nonce}"
    signature = hmac_sha256(server_side_key(), payload)
    register_nonce(actor, nonce)
    store_in_encrypted_session(payload, signature)
    return ok()
\end{lstlisting}

The details are mundane, but the guarantee is valuable: concurrent support sessions remain independently scoped and attributable without shared credentials.
